\chapter{Marco metodológico y gestión del proyecto}
Este capítulo presenta la metodología utilizada para la gestión y desarrollo del proyecto. Se describen las prácticas de 
planificación y seguimiento, la metodología de desarrollo de software adoptada y los instrumentos de recolección de datos 
empleados para identificar las habilidades de los estudiantes. Estos elementos constituyen la base para el diseño, 
implementación y evaluación de la solución propuesta.

\section{Gestión del proyecto (Kanban)}
La organización del proyecto se apoyará principalmente en Kanban, ya que enfoque permite conocer el estado de cada tarea, 
controlar el progreso del proyecto y optimizar la distribución del trabajo. Asimismo, se incorporan prácticas ágiles de 
Scrum, con el propósito de promover un desarrollo ordenado e incremental.

\textbf{Bloques de trabajo:} 
\begin{itemize}
    \item \textbf{Definición del proyecto:} Proporcionará las bases para los requerimientos.
    \item \textbf{Análisis, diseño y arquitectura:} Consolidará su especificación y la solución técnica correspondiente.
    \item \textbf{Desarrollo e implementación:} Construcción e integración de los componentes.
    \item \textbf{Pruebas y resultados:} Evaluación de componentes del sistema.
\end{itemize}

\begin{table}[ht] 
\centering 
\small
\setlength{\tabcolsep}{4pt}
\renewcommand{\arraystretch}{1.15}
\caption{Planeación de bloques de trabajo del proyecto} 
\label{tab:planeacion} 
\begin{tabularx}{\textwidth}{%
    >{\raggedright\arraybackslash}p{0.24\textwidth}
    >{\raggedright\arraybackslash}p{0.18\textwidth}
    >{\raggedright\arraybackslash}X
}
    \toprule 
    \textbf{Bloque de trabajo} & 
    \textbf{Periodo programado} & 
    \textbf{Actividades principales} \\ 
    \midrule 
    Definición del proyecto & 
    8 al 12 de septiembre & 
    Delimitación del problema, definición de objetivos, identificación de la población destinataria, establecimiento del alcance y determinación de necesidades generales. \\ 
    Análisis, diseño y arquitectura & 
    13 al 26 de septiembre & 
    Especificación de requerimientos, diseño de interfaces y estructuras de datos, definición de servicios, contratos de comunicación y procedimientos de análisis. \\ 
    Desarrollo e implementación & 
    27 de septiembre al 19 de octubre & 
    Preparación de datos, desarrollo e integración de servicios, ajuste del modelo de lenguaje, implementación de RAG, entrenamiento y selección del clasificador, construcción de la plataforma y configuración del despliegue. \\
    Pruebas y resultados & 
    20 al 31 de octubre & 
    Ejecución de pruebas, evaluación de modelos y recomendaciones, corrección de incidencias y análisis de resultados. \\ 
    \bottomrule 
\end{tabularx} 
\end{table}

Esta organización permitirá coordinar las actividades necesarias para alcanzar la conclusión del sistema el 31 de octubre y efectuar su 
entrega el 1 de noviembre.

Con el propósito de asegurar el cumplimiento de los objetivos del proyecto, se realizarán verificaciones periódicas durante el 
desarrollo.

\section{Metodología de desarrollo de software (modelo en escalera)}
La metodología de desarrollo de software se basará en el modelo en cascada, el cual organiza el proyecto en etapas secuenciales. 
Para este proyecto se consideran las fases de requerimientos, diseño, implementación y pruebas, cada fase generará entregables que 
servirán como entrada para la siguiente, permitiendo mantener la trazabilidad y el control del proceso de desarrollo. 

\begin{itemize}
    \item \textbf{Requerimientos:} Se establecerán las funciones, reglas, características de calidad, elementos de interfaz y condiciones de implementación que deberá cumplir el sistema. 
    \item \textbf{Diseño:} Se describirán la arquitectura, los modelos de datos, los flujos de información, las interfaces y el procedimiento de procesamiento y clasificación.
    \item \textbf{Implementación:} Se documentará la construcción de los componentes, la preparación de los datos, el desarrollo de los modelos y la integración y el despliegue de la plataforma.
    \item \textbf{Pruebas:} Se presentarán los procedimientos de evaluación, las evidencias obtenidas, las correcciones realizadas y el análisis de los alcances y limitaciones del sistema.
\end{itemize}

\section{Instrumentos de recolección de datos (formularios de habilidades)}


\subsection{Formulario Área 1: Ciencias Físico-Matemáticas e Ingenierías}

\subsection{Formulario Área 2: Ciencias Biológicas, Químicas y de la Salud}

\subsection{Formulario Área 3: Ciencias Sociales y Administrativas}

\endinput
